%% kapa: exam=1 round=20 year=2009 period=1 form=A subject=civil_law subject_ko=민법 seq=35 number=35
%% source: Q-Net 감정평가사 기출문제 (cst00309) artlSeq=5024178 file=extracted/1008/5024178/2009년 제20회 감정평가사 1차시험 기출문제/1교시.hwp
%% conversion: kordoc 4.24.51 -> tools/gen1.py (auto); review=proofread-sonnet; flags=none
\begin{question}{35}
소멸시효와 제척기간의 비교에 관한 설명으로 옳지 \uline{않은} 것은?
\begin{choices}
  \item 소멸시효의 경우 중단사유가 있으면 시효의 진행이 중단되나 제척기간에는 중단제도가 인정되지 아니한다.
  \item 행위무능력을 이유로 한 취소권의 존속기간은 소멸시효기간이 아니라 제척기간이다.
  \item 소멸시효의 완성 후에는 시효이익을 포기할 수 없으나, 제척기간은 기간의 만료로 권리가 당연히 소멸하므로 성질상 포기제도가 없다.
  \item 판례에 따르면 소멸시효의 경우에는 변론주의 원칙상 소멸시효를 주장하여야 참작하나 제척기간의 경우에는 법원이 직권으로 참작하여야 한다.
  \item 소멸시효가 완성되면 권리소멸의 효과가 소급하는데 반하여 제척기간의 경우에는 소급하지 않는다.
\end{choices}
\end{question}
\answer{3}
